\begin{statusdone}
초기 배경 선정과 동적 채우기(사람 박제 방지)가 동작한다.
\end{statusdone}

\chapter{원리 및 설계 근거}

\section{초기 배경 프레임 선정 — \texttt{bg\_select.py}}
동적 채우기에 앞서, 10초간 촬영한 프레임 중 ``기준 배경''을 한 장 고른다.
프레임마다 두 점수를 매긴다.
\begin{itemize}
  \item \textbf{경계 선명도}: 객체 마스크 경계의 색 gradient 평균 — 초점이 맞고
        흔들리지 않은 프레임일수록 높다.
  \item \textbf{레이더 SNR}: 매칭된 객체의 최대 SNR — 레이더가 그 순간 장면을
        잘 보고 있을수록 높다.
\end{itemize}
정규화 후 \[ \text{final} = 0.25\cdot\text{선명도} + 0.75\cdot\text{SNR} \]
로 가중 합산해 최댓값 프레임을 \texttt{background\_raw.jpg}로 저장한다. 이후
ESRGAN 업스케일 → DA3 깊이로 이어진다.

\begin{designbox}{왜 레이더 SNR에 더 큰 가중을 두는가}
선명도만으로 고르면 객체가 많거나 흔들린 순간을 ``경계가 많아'' 선호할 수 있다.
레이더 SNR을 주가중(0.75)으로 두면, 레이더-카메라가 \emph{동시에} 장면을 잘
포착한 프레임이 선택되어 이후 융합 기준으로 더 적합하다.
\end{designbox}

\section{왜 동적 배경인가}
배경 포인트클라우드(3D 뷰어)를 만들려면 객체에 가려지지 않은 깨끗한 배경이
필요하다. 한 장으로는 객체 뒤가 비지만, \emph{시간이 지나며 객체가 비켜난 다른
프레임}에는 그 자리가 드러난다. 여러 프레임에서 ``드러난 배경''을 모아 구멍을
메운다.

\section{누적 채우기 원리}
초기 프레임의 객체 마스크를 ``가려진 영역(remaining)''으로 둔다. 이후 프레임마다
\begin{itemize}
  \item 현재 객체 마스크를 구하고, 각 픽셀이 \emph{연속으로} 객체에 안 잡힌
        횟수(\texttt{reveal\_count})를 센다.
  \item \texttt{REVEAL\_MIN}회 이상 연속으로 비어 있던 픽셀만 ``안정 배경''으로
        인정하고, remaining과 교집합한 영역을 현재 프레임으로 채운다.
  \item 충분한 면적(\texttt{cover\_min})을 새로 채우는 프레임만 선정.
\end{itemize}

\chapter{구현 및 디렉토리 구조}

\section{\texttt{src/background/dynamic\_bg\_fill.py}}
\begin{itemize}
  \item YOLO seg로 매 프레임 객체 마스크(+\texttt{mask\_dilate\_px} 팽창).
  \item \texttt{reveal\_count}로 연속 빈 프레임 카운트, \texttt{REVEAL\_MIN} 이상만 채움.
  \item 누적 채운 양이 \texttt{upscale\_steps}($0.2, 0.25, 0.3$) 단계에 도달하면
        Real-ESRGAN($\times4$, 실패 시 lanczos)으로 업스케일 → \texttt{background.jpg}.
  \item \texttt{\_rerun\_depth(bg)}로 채워진 배경의 깊이를 재추론
        (\texttt{depth\_bg.png}), 뷰어에 재빌드 신호.
\end{itemize}
관련 설정은 \texttt{config/receiver.toml}의 \texttt{[dynamic\_bg]}.

\chapter{문제 해결 과정}

\section{사람 박제(burn-in) 방지}
\textbf{증상}: 가만히 앉은 사람의 상체가 통째로 배경에 굳어 남았다.\\
\textbf{원인}: \texttt{REVEAL\_MIN=3}($\approx0.3$초)으로 너무 짧아, YOLO가 사람을
잠깐 놓친 3프레임을 ``배경''으로 인정하고 박제. 한 번 채우면 remaining에서
빠져 되돌릴 수 없다.\\
\textbf{해결}: \texttt{REVEAL\_MIN}을 15($\approx1.5$초)로 올려 \texttt{receiver.toml}에
노출. 객체가 실제로 1.5초 이상 비켜야 채운다. 마스크 팽창(\texttt{mask\_dilate\_px=15})으로
경계 잔상도 억제.
